\subsection{Corona Performance of the DC Line}
\label{sec:dc-line-corona}

The DC line is not part of the converter supply, but its corona criteria are
set on the same basis as those of the DC switchyard and the station emission
at the DC terminal adds to the noise of the line, so they are given here. On a
DC line the positive pole produces the corona radio noise, and its emission
is highest in fair weather: rain lowers the DC corona noise, the opposite of
the AC case, because the space charge around the positive conductor is
dispersed by the water drops. The design statistic is therefore the fair-weather level not exceeded 50\,\% of the time, at 0.5\,MHz with a quasi-peak
detector, at a lateral distance from the positive pole, typically 15\,m from
the conductor or the edge of the right-of-way. Values of 40\,dB($\mu$V/m) to
50\,dB($\mu$V/m) at the edge of the right-of-way are the range of the design
criteria used for $\pm$500\,kV lines in the western United States~[37], chosen
to protect AM broadcast reception at the signal strengths found along the
routes. The matching audible-noise criterion is an A-weighted level of
42\,dB(A) in fair weather, 50\,\% statistic, at the same point. Both are
calculated for the maximum elevation of the route; a route through Utah and
Nevada runs between 1200\,m and 2000\,m, where the relative air density lies
between 0.87 and 0.79 and the corona onset gradient falls by the same ratio.
\FloatBarrier
